\documentclass[11pt]{article}
\usepackage{amsmath,amssymb,amsthm,algorithm,algorithmic}
\usepackage{booktabs}
\usepackage[margin=1in]{geometry}

\newtheorem{theorem}{Theorem}
\newtheorem{lemma}[theorem]{Lemma}
\newtheorem{corollary}[theorem]{Corollary}
\newtheorem{proposition}[theorem]{Proposition}
\newtheorem{definition}[theorem]{Definition}

\title{Problem 3: Largest Prime Factor}
\author{Project Euler Solutions}
\date{}

\begin{document}
\maketitle

\section{Problem Statement}

What is the largest prime factor of the number $600\,851\,475\,143$?

\section{Mathematical Development}

\subsection{Fundamental Theorems}

\begin{definition}
An integer $p \ge 2$ is \emph{prime} if its only positive divisors are $1$ and $p$. An integer $n \ge 2$ that is not prime is \emph{composite}.
\end{definition}

\begin{theorem}[Smallest Divisor is Prime]
\label{thm:smallest}
If $n \ge 2$ is an integer and $d$ is the smallest integer $\ge 2$ that divides $n$, then $d$ is prime.
\end{theorem}

\begin{proof}
Suppose for contradiction that $d$ is composite: $d = ab$ with $1 < a < d$ and $b > 1$. Since $a \mid d$ and $d \mid n$, transitivity gives $a \mid n$ with $2 \le a < d$, contradicting the minimality of~$d$.
\end{proof}

\begin{theorem}[Square Root Bound]
\label{thm:sqrt}
If $n \ge 2$ is composite, then $n$ has a prime factor $p \le \lfloor\sqrt{n}\rfloor$.
\end{theorem}

\begin{proof}
Write $n = ab$ with $1 < a \le b < n$. Then $a^2 \le ab = n$, so $a \le \sqrt{n}$. Let $p$ be the smallest divisor of $a$ with $p \ge 2$. By Theorem~\ref{thm:smallest}, $p$ is prime. Since $p \mid a$ and $a \mid n$, we have $p \mid n$ with $p \le a \le \sqrt{n}$.
\end{proof}

\begin{corollary}[Primality Test]
\label{cor:prime-test}
An integer $n \ge 2$ is prime if and only if no prime $p \le \lfloor\sqrt{n}\rfloor$ divides $n$.
\end{corollary}

\begin{proof}
The forward direction is trivial. The contrapositive of the reverse direction is Theorem~\ref{thm:sqrt}.
\end{proof}

\subsection{Correctness of Trial Division}

\begin{theorem}[Correctness]
\label{thm:correct}
The trial division procedure (Algorithm~1) correctly produces the prime factorization of $n$ and identifies the largest prime factor.
\end{theorem}

\begin{proof}
Let $n_0$ denote the original input. We maintain the invariant:

\smallskip
\noindent\textbf{Invariant $\mathcal{I}(d)$:} $n_0 = n \cdot Q$ where $Q$ is the product of all extracted factors, and $n$ has no prime factor $< d$.
\smallskip

\emph{Initialization.} At $d=2$: $Q=1$, $n = n_0$, and ``no prime factor $< 2$'' is vacuous.

\emph{Maintenance.} When $d \mid n$: by $\mathcal{I}(d)$, the smallest divisor of $n$ that is $\ge 2$ is $\ge d$; since $d \mid n$, this smallest divisor is exactly $d$. By Theorem~\ref{thm:smallest}, $d$ is prime. We divide out all copies of $d$, then increment $d$, preserving $\mathcal{I}(d+1)$.

\emph{Termination.} When $d^2 > n$ and $n > 1$: if $n$ were composite, Theorem~\ref{thm:sqrt} would give a prime factor $\le \sqrt{n} < d$, contradicting $\mathcal{I}(d)$. Hence $n$ is prime and is the largest prime factor.
\end{proof}

\begin{theorem}[Termination]
\label{thm:term}
Trial division terminates for all inputs $n \ge 2$.
\end{theorem}

\begin{proof}
In each inner iteration $n$ decreases by a factor $\ge 2$; in each outer iteration $d$ increases by~1. The measure $\mu = n + (\lfloor\sqrt{n}\rfloor - d)$ strictly decreases, and the loop condition ensures $\mu \ge 0$.
\end{proof}

\subsection{Factorization of $600\,851\,475\,143$}

\begin{lemma}
\label{lem:factorization}
$600\,851\,475\,143 = 71 \times 839 \times 1471 \times 6857$.
\end{lemma}

\begin{proof}
Direct computation: $71 \times 839 = 59\,569$; then $59\,569 \times 1471 = 87\,625\,999$; then $87\,625\,999 \times 6857 = 600\,851\,475\,143$.
\end{proof}

\begin{lemma}
\label{lem:primes}
Each of $71$, $839$, $1471$, $6857$ is prime.
\end{lemma}

\begin{proof}
By Corollary~\ref{cor:prime-test}, it suffices to verify that no prime $p \le \lfloor\sqrt{n}\rfloor$ divides each factor $n$:
\begin{itemize}
    \item $71$: $\lfloor\sqrt{71}\rfloor = 8$. Check primes $2,3,5,7$: $71$ is odd, $3 \nmid 71$ (digit sum $8$, $3\nmid 8$), not divisible by 5 or 7.
    \item $839$: $\lfloor\sqrt{839}\rfloor = 28$. Check primes $2,3,5,7,11,13,17,19,23$: none divide $839$.
    \item $1471$: $\lfloor\sqrt{1471}\rfloor = 38$. Check primes up to $37$: none divide $1471$.
    \item $6857$: $\lfloor\sqrt{6857}\rfloor = 82$. Check primes up to $79$: none divide $6857$.
\end{itemize}
\end{proof}

\begin{theorem}
The largest prime factor of $600\,851\,475\,143$ is $6857$.
\end{theorem}

\begin{proof}
By Lemma~\ref{lem:factorization}, $600\,851\,475\,143 = 71 \times 839 \times 1471 \times 6857$. By Lemma~\ref{lem:primes}, all factors are prime. By the Fundamental Theorem of Arithmetic, this factorization is unique. The largest factor is $6857$.
\end{proof}

\section{Algorithm}

\begin{algorithm}[H]
\caption{\textsc{LargestPrimeFactor}$(n)$}
\begin{algorithmic}[1]
\STATE $d \gets 2$
\WHILE{$d \cdot d \le n$}
    \WHILE{$n \bmod d = 0$}
        \STATE $n \gets n / d$
    \ENDWHILE
    \STATE $d \gets d + 1$
\ENDWHILE
\IF{$n > 1$}
    \RETURN $n$
\ELSE
    \RETURN $d - 1$
\ENDIF
\end{algorithmic}
\end{algorithm}

\textbf{Execution trace for $n = 600\,851\,475\,143$:}

\begin{center}
\begin{tabular}{rll}
\toprule
$d$ & Action & Remaining $n$ \\
\midrule
$71$   & $71 \mid n$; divide & $8\,462\,696\,833$ \\
$839$  & $839 \mid n$; divide & $10\,086\,647$ \\
$1471$ & $1471 \mid n$; divide & $6\,857$ \\
$1472$ & $d^2 > n$; exit & $6\,857$ \\
\bottomrule
\end{tabular}
\end{center}

Return $n = 6857$.

\section{Complexity Analysis}

\begin{theorem}
Trial division runs in $O(\sqrt{n_0})$ time and $O(1)$ space, where $n_0$ is the input.
\end{theorem}

\begin{proof}
\emph{Time.} The outer loop increments $d$ from $2$ up to at most $\lfloor\sqrt{n_0}\rfloor$: at most $\sqrt{n_0}$ outer iterations. The total number of inner iterations across all $d$ is at most $\lfloor\log_2 n_0\rfloor$ (each halves $n$ at minimum). Total: $O(\sqrt{n_0} + \log n_0) = O(\sqrt{n_0})$.

\emph{Tightness.} When $n_0$ is prime, the inner loop never executes and the outer loop runs $\lfloor\sqrt{n_0}\rfloor - 1$ times, giving $\Theta(\sqrt{n_0})$.

\emph{Space.} Only the variables $d$ and $n$ are stored: $O(1)$.

For $n_0 = 600\,851\,475\,143$, the worst-case bound is $\sqrt{n_0} \approx 775\,146$ iterations.
\end{proof}

\section{Answer}

\[
\boxed{6857}
\]

\end{document}
